\section{作文素材}

\subsection{高频主题素材库}

\subsubsection{主题一：时代担当与青年责任}

\begin{itemize}
  \item \textbf{人物素材}：黄文秀（脱贫攻坚）、陈祥榕（戍边英雄）、张桂梅（教育扶贫）
  \item \textbf{时事素材}：中国空间站建成、一带一路十周年、乡村振兴战略
  \item \textbf{名言}：
    \begin{itemize}
      \item 「青年之字典，无『困难』之字；青年之口头，无『障碍』之语。」—— 李大钊
      \item 「为天地立心，为生民立命，为往圣继绝学，为万世开太平。」—— 张载
      \item 「一个时代的精神是青年代表的精神。」—— 马克思
    \end{itemize}
  \item \textbf{运用角度}：个人理想与国家命运的结合；平凡岗位的不凡贡献
\end{itemize}

\subsubsection{主题二：科技发展与人文关怀}

\begin{itemize}
  \item \textbf{人物素材}：屠呦呦（青蒿素挽救百万生命）、南仁东（FAST 天眼之父）
  \item \textbf{时事素材}：ChatGPT 引发教育思考、人脸识别与隐私边界
  \item \textbf{名言}：
    \begin{itemize}
      \item 「科学没有国界，但科学家有祖国。」—— 巴斯德
      \item 「科技是一种手段，人文才是目的。」—— 爱因斯坦
      \item 「我们往往只关心如何走得更快，却忘了为什么出发。」
    \end{itemize}
  \item \textbf{运用角度}：科技的双刃剑；工具理性与价值理性的平衡
\end{itemize}

\subsubsection{主题三：文化传承与创新发展}

\begin{itemize}
  \item \textbf{人物素材}：樊锦诗（敦煌守护人）、单霁翔（故宫活化）、李子柒（文化输出）
  \item \textbf{时事素材}：国潮兴起（汉服、文创）、非遗传承人培养计划
  \item \textbf{名言}：
    \begin{itemize}
      \item 「只有民族的，才是世界的。」—— 鲁迅
      \item 「文化兴国运兴，文化强民族强。」
      \item 「守正不守旧，尊古不复古。」
    \end{itemize}
  \item \textbf{运用角度}：传统与现代的融合；文化自信的表现；从文化自觉到文化自信
\end{itemize}

\subsubsection{主题四：个人品格与精神追求}

\begin{itemize}
  \item \textbf{人物素材}：袁隆平（脚踏实地）、钟南山（逆行而上）、司马迁（坚韧不拔）
  \item \textbf{时事素材}：北京冬奥会运动员拼搏精神、航天员矢志报国
  \item \textbf{名言}：
    \begin{itemize}
      \item 「天行健，君子以自强不息。」——《周易》
      \item 「志之所趋，无远弗届；穷山距海，不能限也。」——《格言联璧》
      \item 「千磨万击还坚劲，任尔东西南北风。」—— 郑板桥
    \end{itemize}
  \item \textbf{运用角度}：面对挫折的态度；坚守与变通；平凡与伟大的辩证
\end{itemize}

\subsubsection{主题五：合作共赢与命运与共}

\begin{itemize}
  \item \textbf{人物素材}：中国援外医疗队、联合国维和部队中的中国力量
  \item \textbf{时事素材}：「一带一路」国际合作、气候变化的全球治理
  \item \textbf{名言}：
    \begin{itemize}
      \item 「一花独放不是春，百花齐放春满园。」
      \item 「人类生活在同一个地球村里。」—— 习近平
      \item 「各美其美，美人之美，美美与共，天下大同。」—— 费孝通
    \end{itemize}
  \item \textbf{运用角度}：竞争与合作的关系；全球化背景下的团结
\end{itemize}

\subsection{综合人物素材库（按时代分类）}

\subsubsection{古代人物}

\begin{center}
\begin{tabular}{c|c|c}
\textbf{人物} & \textbf{一句话素材} & \textbf{可用主题} \\ \hline
屈原 & 「路漫漫其修远兮，吾将上下而求索。」——在放逐中写下《离骚》，以生命捍卫理想 & 爱国、理想、坚守、不屈 \\
司马迁 & 遭宫刑而著《史记》——「人固有一死，死有重于泰山，或轻于鸿毛」 & 坚韧、使命、信念、逆境 \\
诸葛亮 & 「鞠躬尽瘁，死而后已。」——出师一表真名世，千载谁堪伯仲间 & 忠诚、担当、智慧、奉献 \\
范仲淹 & 「先天下之忧而忧，后天下之乐而乐。」——不以物喜不以己悲的士大夫精神 & 家国情怀、担当、格局 \\
苏轼 & 三度被贬，却写下「一蓑烟雨任平生」「此心安处是吾乡」 & 豁达、逆境、乐观、文化自信 \\
文天祥 & 「人生自古谁无死，留取丹心照汗青。」——从容就义，气节千古 & 爱国、气节、生死观、信仰 \\
王阳明 & 龙场悟道，提出「知行合一」——事上磨练，方得真知 & 实践、知行、智慧、修身 \\
林则徐 & 「苟利国家生死以，岂因祸福避趋之。」——虎门销烟，开眼看世界 & 爱国、担当、改革、勇气 \\
\end{tabular}
\end{center}

\subsubsection{近现代人物}

\begin{center}
\begin{tabular}{c|c|c}
\textbf{人物} & \textbf{一句话素材} & \textbf{可用主题} \\ \hline
鲁迅 & 「横眉冷对千夫指，俯首甘为孺子牛。」——以笔为刀，唤醒国人 & 批判精神、担当、独立思考 \\
钱学森 & 放弃美国优厚待遇回国——「我的事业在中国，我的成就在中国」 & 爱国、奉献、科学精神、选择 \\
邓稼先 & 隐姓埋名 28 年研制原子弹——「不要让人家把我们落得太远」 & 奉献、爱国、坚守、牺牲 \\
于敏 & 「国产专家一号」——从零开始突破氢弹，填补中国核武器空白 & 自主创新、科学报国、淡泊名利 \\
袁隆平 & 把一生浸在稻田里——「我不在试验田，就在去试验田的路上」 & 实干、奉献、科学精神、民生 \\
张桂梅 & 创办免费女高，把 1800 多个女孩送出大山——「只要还有一口气，就要站在讲台上」 & 奉献、教育、信念、女性力量 \\
\end{tabular}
\end{center}

\subsubsection{当代青年榜样}

\begin{center}
\begin{tabular}{c|c|c}
\textbf{人物} & \textbf{一句话素材} & \textbf{可用主题} \\ \hline
黄文秀 & 北师大硕士回乡扶贫，驻村里程两万五千里——用生命走完新时代长征路 & 青春奉献、理想、担当 \\
陈祥榕 & 19 岁戍边牺牲——「清澈的爱，只为中国」 & 爱国、青春、奉献、担当 \\
苏炳添 & 32 岁跑出 9 秒 83——「我想证明，黄种人也可以飞」 & 突破极限、坚持、爱国、自律 \\
全红婵 & 14 岁奥运冠军，来自贫困家庭——「妈妈治病要花很多钱，我想赚钱给她治病」 & 拼搏、孝心、纯真、奋斗 \\
江梦南 & 半聋哑女孩考上清华博士——「我从来没有因为听不见，就把自己看成一个弱者」 & 自强、逆境、教育、信念 \\
陈贝儿 & 香港记者拍摄《无穷之路》——用镜头讲述中国脱贫故事 & 文化传播、真实记录、家国认同 \\
航天员群体 & 「感觉良好」背后的千锤百炼——平均年龄不到 40 岁的中国航天 & 团队、奉献、科技报国、青年 \\
\end{tabular}
\end{center}

\subsubsection{名言分类速查表}

\begin{center}
\begin{tabular}{c|c}
\textbf{主题} & \textbf{可用名言} \\ \hline
爱国 & «位卑未敢忘忧国»——陆游；«天下兴亡，匹夫有责»——顾炎武 \\
理想 & «志当存高远»——诸葛亮；«心之所向，素履以往» \\
坚持 & «千淘万漉虽辛苦，吹尽狂沙始到金»——刘禹锡 \\
创新 & «苟日新，日日新，又日新»——《大学》；«周虽旧邦，其命维新» \\
奋斗 & «青春须早为，岂能长少年»——孟郊；«幸福都是奋斗出来的» \\
奉献 & «春蚕到死丝方尽，蜡炬成灰泪始干»——李商隐 \\
自信 & «天生我材必有用»——李白；«自信人生二百年，会当水击三千里»——毛泽东 \\
格局 & «牢骚太盛防肠断，风物长宜放眼量»——毛泽东 \\
实践 & «纸上得来终觉浅，绝知此事要躬行»——陆游 \\
惜时 & «少年易老学难成，一寸光阴不可轻»——朱熹 \\
\end{tabular}
\end{center}

\subsection{素材积累引导：从「背素材」到「用素材」}

\subsubsection{素材使用三原则}

\begin{enumerate}
  \item \textbf{精准不贪多}：考场上一篇文章用好 2--3 个素材即可，不要堆砌
  \item \textbf{有细节才鲜活}：空洞地提「袁隆平很伟大」不如写「袁隆平 90 岁每天下田」
  \item \textbf{分析重于罗列}：写完素材必须分析「这个事例说明了什么」，扣回论点
\end{enumerate}

\subsubsection{素材运用示范：同一个素材在不同位置的使用}

\textbf{以「樊锦诗」为例，展示一个素材如何贯穿全文}：

\begin{itemize}
  \item[\textbf{开头段}] 引用式引入：「从北大才女到敦煌守护人，樊锦诗用一生诠释了什么叫『择一事，终一生』。」
  \item[\textbf{主体段分论点}]「真正的热爱，经得起时间的打磨。」樊锦诗初到敦煌时，住土房、点油灯、喝咸水，很多人待不了三年就走了。她却一待就是五十多年，从青丝到白发。这不是苦行，而是因为她看见了敦煌的美——那种穿越千年的文明之美，足以让一个人甘愿付出一生去守护。（PEE 结构完整演示）
  \item[\textbf{结尾段}] 升华：「樊锦诗在敦煌找到了一生所爱，而我们这一代青年，又该在哪里安放我们的热爱？时代给了我们更广阔的舞台——无论是乡村振兴的田野、科技攻关的实验室，还是文化传承的第一线，都值得我们去『择一事，终一生』。」
\end{itemize}

\subsubsection{万金油主题转化表——任何题目都能落脚爱国/青年奉献}

\begin{center}
\begin{tabular}{c|c|c}
\textbf{考题话题} & \textbf{转化逻辑链条} & \textbf{落脚点} \\ \hline
科技/AI & 科技需要方向 $\to$ 谁来掌舵？$\to$ 有家国情怀的青年 & 青年以科技报国 \\
创新/守正 & 创新的目的是什么？$\to$ 为国家发展注入动力 & 青年勇担创新使命 \\
文化传承 & 为什么要传承？$\to$ 文化自信是民族复兴之基 & 青年做文化传承者 \\
个人理想 & 理想的价值在于为谁而立志 $\to$ 融入时代洪流 & 青年把个人理想融入国家需要 \\
挫折/困难 & 面对困难的态度 $\to$ 时代需要坚韧的青年 & 青年在奋斗中成长 \\
合作/团结 & 团结的力量 $\to$ 集中力量办大事的制度优势 & 青年投身国家大局 \\
阅读/学习 & 学以致用 $\to$ 为谁而学？$\to$ 为民族复兴而读书 & 青年以学报国 \\
规则/自律 & 规矩意识 $\to$ 社会需要怎样的青年 & 青年做时代新风的引领者 \\
生态/环保 & 绿水青山就是金山银山 $\to$ 代代守护 & 青年守护美丽中国 \\
平凡/伟大 & 平凡岗位的不凡贡献 $\to$ 每个人都是时代的建设者 & 青年在平凡中铸就伟大 \\
\end{tabular}
\end{center}

\textbf{使用示范}：假如考题是「科技与人文」，直接查表找到第一行：

\begin{enumerate}
  \item 开头正常论述科技与人文的关系
  \item 分论点 1--2 围绕话题本身展开
  \item 分论点 3 自然递进：「科技向善需要方向，而方向的确立需要一代有理想、有担当的青年。当代青年应当以家国情怀为科技创新的导航仪，用智慧与汗水在科技强国的征程上写下属于我们的答卷。」
\end{enumerate}

关键词切换技巧——用这些词把话题「拧」向爱国/奉献：

\[
\textbf{「时代」「使命」「担当」「征程」「答卷」「复兴」「家园」}
\]

每个主体段或结尾段出现 1--2 个这类词汇，整篇文章的格调自然提升，且完全不会显得生硬。

\subsection{素材转化为论述文段——手把手拆解}

很多同学背了一堆素材，写出来却是「袁隆平很伟大，我们要向他学习」——空洞、没有逻辑。问题出在：\textbf{不知道素材和论点之间缺了一座桥}。

这座桥就是\textbf{论证}。下面从最基础的步骤开始，一步步拆解如何把一段素材变成一篇严谨的论述段落。

\subsubsection{第一步：给素材「找角度」——同一个素材，多个切面}

拿到一个素材，先别急着背，问自己三个问题：

\begin{enumerate}
  \item 这个人做了什么？（事实概括）
  \item 他为什么这么做？（动机分析）
  \item 这反映了他什么品质/精神？（提炼价值）
\end{enumerate}

\textbf{以「袁隆平」为例，拆出 6 个可用角度}：

\begin{center}
\begin{tabular}{c|c|c|c}
\textbf{角度} & \textbf{素材细节} & \textbf{提炼关键词} & \textbf{适合的作文话题} \\ \hline
1 & 90 岁还每天下田 & 坚守、实干 & 坚持、脚踏实地、敬业 \\
2 & 让中国人把饭碗端在自己手里 & 爱国、粮食安全 & 爱国、科技报国 \\
3 & 在杂交水稻被质疑时坚持研究 & 信念、定力 & 逆境、坚持自我 \\
4 & 把野生稻基因保留下来 & 前瞻、传承 & 长远眼光、基础研究 \\
5 & 无偿将杂交水稻技术分享给世界 & 格局、人类关怀 & 合作共赢、大国担当 \\
6 & 一生只做一件事 & 专注、匠心 & 专注、工匠精神 \\
\end{tabular}
\end{center}

\textbf{练习}：自己试着给「张桂梅」找出 5 个以上的角度（参考答案在文末）。

\subsubsection{第二步：从「角度」到「论点句」——把角度和考题连接起来}

有了角度还不够，要把角度和 \textbf{考题的关键词} 焊接在一起。

\textbf{公式}：
\[
\text{角度关键词} + \text{考题关键词} + \text{判断（是/不是/应该/不应该）} = \text{论点句}
\]

\textbf{示范}（假设考题是「坚守与创新」）：

\begin{itemize}
  \item[素材] 袁隆平 90 岁下田（角度：坚守）
  \item[论点句] 真正的创新，从来不是灵光一现，而是建立在数十年如一日坚守之上的厚积薄发。
\end{itemize}

\textbf{示范 2}（假设考题是「个人价值与社会贡献」）：

\begin{itemize}
  \item[素材] 袁隆平分享杂交水稻技术（角度：格局）
  \item[论点句] 一个人的价值，不在于他从世界获取了多少，而在于他为世界贡献了什么。
\end{itemize}

\textbf{检验标准}：去掉论点句中的素材人名，如果读者能猜到你在写什么话题——说明论点句合格了。

\subsubsection{第三步：从「论点句」到「论述段落」——PEE 完整展开}

有了论点句，接下来把它展开成一个完整段落。

\textbf{PEE 公式回顾}：

\[
\boxed{\text{Point 提出分论点}} \to \boxed{\text{Evidence 叙述素材}} \to \boxed{\text{Explanation 分析论证}}
\]

\textbf{注意}：PEE 的难点不在 P 和 E，在 \textbf{第二个 E（分析）}。很多同学只写到 E 就停了——「袁隆平 90 岁下田，所以他很伟大」——这是\textbf{陈述}，不是\textbf{论证}。

分析的本质是回答一句话：\textbf{「这个例子为什么能证明我的观点？」}

\subsubsection{分析的五种基本逻辑}

\begin{center}
\begin{tabular}{c|c|c}
\textbf{逻辑类型} & \textbf{分析句式} & \textbf{示例} \\ \hline
因果分析 & 「之所以……是因为……」 & 袁隆平之所以 90 岁还在下田，不是因为他闲不住，而是因为他把「让所有人吃饱饭」刻进了生命的本能里。 \\
假设分析 & 「试想，如果……那么……」 & 试想，如果袁隆平在取得初步成功后选择退休享福，杂交水稻还能一次次突破亩产极限吗？ \\
对比分析 & 「与……不同，……」 & 与那些追求短期利益的人不同，袁隆平选择了一条最难的路——用几十年时间做一件不一定能看到结果的事。 \\
本质分析 & 「……的本质不是……而是……」 & 袁隆平 90 岁下田，表面上看是一种习惯，本质上是一种信仰——他相信土地不会辜负每一个认真对待它的人。 \\
升华分析 & 「……不仅……更……」 & 袁隆平留下的不仅是杂交水稻，更是一种精神——一种把个人命运与国家和人类命运紧紧绑在一起的精神。 \\
\end{tabular}
\end{center}

\textbf{核心原则}：\textbf{每一种分析，最后都要落回到分论点句上}。写完分析句，读一遍，看它和段落第一句（Point）是否有逻辑关联。如果没有，说明分析偏了。

\subsubsection{完整拆解示范 1：从素材到段落的一步一讲}

\textbf{任务}：以「黄文秀」为素材，写一段论证「青春的价值在于选择」。

\begin{center}
\begin{tabular}{c|l|l}
\textbf{步骤} & \textbf{内容} & \textbf{讲解} \\ \hline
素材 & 黄文秀，北师大硕士，放弃城市工作机会，回到广西百坭村当第一书记。驻村满一年时，她的汽车里程达到两万五千公里。2019 年暴雨之夜，她驱车回村途中遭遇山洪，30 岁牺牲。 & 先明确素材是什么 \\
角度 & 选择——她放弃了城市选择了农村 & 从素材中提取可用角度 \\
P 论点句 & 青春的价值，不在于选择了什么，而在于为什么而选择。 & 论点句必须出现考题关键词 \\
E 叙事 & 黄文秀从北师大硕士毕业后，没有像大多数人一样留在大城市，而是回到了广西最贫困的百坭村。她驻村一年的行驶里程，恰好是红军长征的长度——两万五千里。 & 叙事要有细节（两万五千里），不要流水账 \\
E 因果分析 & 她为什么做出这样的选择？因为她深知，北师大的学历不是向上爬的阶梯，而是向下扎根的底气。一个人的价值，不在于占据什么位置，而在于为多少人带来了改变。 & 用因果分析回答「为什么」\\
回扣 & 黄文秀用生命告诉我们：青春最珍贵的，不是选择了安逸，而是选择了担当。 & 最后一句话必须扣回 P 中的关键词 \\
\end{tabular}
\end{center}

\textbf{完整段落}（考场直接写出来的样子）：

「青春的价值，不在于选择了什么，而在于为什么而选择。（P）黄文秀从北师大硕士毕业后，没有像大多数人一样留在大城市，而是回到了广西最贫困的百坭村。她驻村一年的行驶里程，恰好是红军长征的长度——两万五千里。（E）她为什么做出这样的选择？因为她深知，北师大的学历不是向上爬的阶梯，而是向下扎根的底气。一个人的价值，不在于占据什么位置，而在于为多少人带来了改变。（E）黄文秀用生命告诉我们：青春最珍贵的，不是选择了安逸，而是选择了担当。（回扣）」

\subsubsection{完整拆解示范 2：从素材到段落——对比论证}

\textbf{任务}：以「苏轼」为素材，写一段论证「困境中如何保持精神自由」。

\begin{center}
\begin{tabular}{c|l|l}
\textbf{步骤} & \textbf{内容} & \textbf{讲解} \\ \hline
素材 & 苏轼一生三度被贬：黄州、惠州、儋州，一次比一次偏远。但他每到一处都活得精彩——在黄州写《赤壁赋》，在惠州发明荔枝酒，在儋州办学教书。 & 选素材中最能体现主题的部分 \\
P 论点句 & 真正的精神自由，不是环境给你的，而是你自己修来的。 & 论点要鲜明，有态度 \\
E 叙事 & 苏轼一生三度被贬，从黄州到惠州再到儋州，一次比一次荒远。换了别人，可能早已被命运击垮。但苏轼在黄州写下了「大江东去」，在惠州发明了荔枝酒，在儋州开办学堂，把文明的种子播撒到天涯海角。 & 叙事中用对比（「换了别人」） \\
E 对比分析 & 与那些被困境压垮的人不同，苏轼的可贵在于：他从不把命运的波折当作堕落的借口。被贬不是终点，而是另一种生活的起点。因为他心中装着一个足够大的世界——有江山风月，有诗酒文章，有黎民百姓——所以任何牢笼都关不住他。 & 对比分析 + 本质分析 \\
回扣 & 苏轼用一生证明：只要精神不被囚禁，就没有什么能真正打败你。 & 回扣 P 中的「精神自由」 \\
\end{tabular}
\end{center}

\subsubsection{完整拆解示范 3：从素材到段落——层进式论证（最难也最高分）}

\textbf{任务}：以「屠呦呦」为素材，写一段论证「创新需要什么」。

\begin{center}
\begin{tabular}{c|l|l}
\textbf{步骤} & \textbf{内容} & \textbf{讲解} \\ \hline
P 论点句 & 真正的创新，从来不是凭空产生，而是扎根传统、拥抱失败、心系苍生的结果。 & 论点句本身就包含了三个层次 \\
E 分层 1 & 扎根传统：屠呦呦翻阅大量古籍，从葛洪《肘后备急方》中「青蒿一握，以水二升渍，绞取汁，尽服之」获得了关键灵感。 & 第一层：创新需要继承 \\
E 分层 2 & 拥抱失败：在 190 次失败后，她带领团队重新审视实验流程——不是药材错了，而是提取方法错了。高温破坏了青蒿的有效成分。 & 第二层：创新需要试错 \\
E 分层 3 & 心系苍生：她亲自参与临床试验，不顾自身健康。因为在她看来，科研的目的不是发表论文，而是救人。 & 第三层：创新需要目的 \\
回扣 & 屠呦呦的故事告诉我们：最好的创新，是站在前人的肩膀上、试过千百条错路后，带着对生命的敬畏，找到那一株改变世界的青蒿。 & 综合回扣三个层次 \\
\end{tabular}
\end{center}

\subsubsection{第四步：常见逻辑错误自查清单}

写完一个段落，用以下 6 个问题逐个检查：

\begin{enumerate}
  \item \textbf{论点句是否清晰}？——整段第一句话，能不能让读者一眼看到你的观点？
  \item \textbf{素材和论点是否匹配}？——这个例子真的能证明你的观点吗？（常见错误：用勤奋的例子去证明创新的观点）
  \item \textbf{有没有分析句}？——不能叙事完就结束，必须有「这说明……」之类的分析句
  \item \textbf{分析是否扣题}？——分析的内容必须回扣到段落的第一句（P），不能偏到另一个话题
  \item \textbf{有没有过渡词}？——「之所以……是因为……」「试想……」「与……不同」——这些词让逻辑关系显性化
  \item \textbf{字数是否合适}？——一段 180--220 字，P 约 20 字 + E 约 60 字 + E（分析）约 100 字 + 回扣约 20 字
\end{enumerate}

\subsubsection{从段落拼成整篇文章——3 段 PEE 如何组合}

\begin{center}
\begin{tabular}{c|c|c}
\textbf{位置} & \textbf{功能} & \textbf{论述方式建议} \\ \hline
分论点 1 & 个人层面——话题对我的启发 & 因果分析为主 \\
分论点 2 & 社会层面——话题对社会的影响 & 对比分析为主 \\
分论点 3 & 国家/时代层面——话题与青年使命 & 升华分析为主 \\
\end{tabular}
\end{center}

每个分论点段都用 PEE 结构，但三个段落的 \textbf{分析方式要有所变化}，避免单调。分论点 3 用升华分析自然过渡到爱国/青年奉献，正是之前学过的万金油收束法。

\textbf{完整示例}：把以上三个示范段落（黄文秀、苏轼、屠呦呦）拼入同一篇文章的框架：

\begin{itemize}
  \item[\textbf{开头}] 引出话题 + 中心论点
  \item[\textbf{分论点 1}] 个人成长需要正确选择（用黄文秀素材——因果分析）
  \item[\textbf{分论点 2}] 面对困境需要精神自由（用苏轼素材——对比分析）
  \item[\textbf{分论点 3}] 时代进步需要创新担当（用屠呦呦素材——升华分析 $\to$ 落脚青年奉献）
  \item[\textbf{结尾}] 总结 + 升华
\end{itemize}

\subsubsection{「张桂梅」角度拆解练习答案参考}

\begin{center}
\begin{tabular}{c|c|c}
\textbf{角度} & \textbf{素材细节} & \textbf{可写话题} \\ \hline
1 & 创办全国第一所免费女子高中 & 教育公平、女性力量 \\
2 & 身患 23 种疾病仍坚持早起陪学生晨读 & 奉献、坚守、师者仁心 \\
3 & 把 1800 多个女孩送出大山 & 改变命运、希望 \\
4 & 12 年家访行程 11 万公里 & 脚踏实地、信念 \\
5 & 「我生来就是高山而非溪流」——华坪女高校训 & 自信、突破偏见、女性自立 \\
6 & 被评为「七一勋章」获得者 & 奉献被看见、荣誉与责任 \\
\end{tabular}
\end{center}

\subsection{万能开头结尾}

\subsubsection{万能开头句式}

\begin{enumerate}
  \item \textbf{引用式}：「……曾说：『……』。这句话启示我们，……。在当今时代，……更显其深刻意义。」
  \item \textbf{设问式}：「面对……，我们不禁要问：……？答案或许并不唯一，但一个清晰的指向是——……。」
  \item \textbf{排比式}：「……，是……；……，是……；……，是……。而……，正是这一切的基石。」
  \item \textbf{对比式}：「古有……，今有……。时代的车轮滚滚向前，但……的价值从未改变。」
  \item \textbf{现象式}：「在……的时代背景下，……的现象日益凸显。如何……，成为我们无法回避的问题。」
\end{enumerate}

\subsubsection{万能结尾句式}

\begin{enumerate}
  \item \textbf{总结升华式}：「总之，……不仅是……，更是……。作为新时代青年，我们应……，以……的姿态，书写属于我们这一代人的答卷。」
  \item \textbf{呼吁号召式}：「让我们……，从自身做起，从现在做起。唯其如此，……才能真正成为现实。」
  \item \textbf{展望未来式}：「……之路并非坦途，但我们坚信……。待到……之时，便是……之日。」
  \item \textbf{首尾呼应式}：「开篇所问……，至此已有了明确的答案：……。诚如……所言……，这正是……的真谛。」
\end{enumerate}

\subsection{优秀范文示例}

\subsubsection{范文 1：时代担当类}

\textbf{题目}：《以青春之我，担时代之责》

\textbf{精彩开头}：
「每一代青年都有自己的际遇和机缘，都要在自己所处的时代条件下谋划人生、创造历史。从五四运动的街头呐喊，到改革开放的创业浪潮，再到新时代的星辰大海，中国青年始终与民族复兴同频共振。」

\textbf{主体结构}：
\begin{itemize}
  \item 分论点 1：担当需要「功成不必在我」的境界
  \item 分论点 2：担当需要「功成必定有我」的决心
  \item 分论点 3：担当需要「踏平坎坷成大道」的实干
\end{itemize}

\textbf{精彩结尾}：
「历史的接力棒已经传递到我们手中。我们是什么样子，未来的中国就是什么样子。以青春之我，担时代之责，这是我们对历史最好的回答，也是我们对未来最庄严的承诺。」

\subsubsection{范文 2：思辨关系类}

\textbf{题目}：《在传承中创新，在创新中传承》

\textbf{精彩开头}：
「有人说传承是守旧，有人说创新是忘本。其实，传承与创新并非对立的两极，而是一体两面。没有传承的创新是无源之水，没有创新的传承是死水一潭。」

\textbf{主体结构}：
\begin{itemize}
  \item 分论点 1：传承是根基，让文化有魂
  \item 分论点 2：创新是动力，让文化有生命力
  \item 分论点 3：传承与创新共生，方能让文明生生不息
\end{itemize}

\textbf{精彩结尾}：
「五千年的文明告诉我们：守正不守旧，尊古不复古。在传承中创新，在创新中传承——这才是文化自信的真正内涵，也是我们这一代人对中华文明最好的守护。」

\subsubsection{范文 3：个人成长类}

\textbf{题目}：《在快时代里慢下来》

\textbf{写作特点}：
\begin{itemize}
  \item 切入点小：从「倍速追剧」「碎片阅读」的生活现象入手
  \item 层层递进：现象分析 $\to$ 原因探究 $\to$ 解决方案
  \item 结尾有力：「慢下来不是停滞，而是为了更好地出发。」
\end{itemize}

\subsubsection{范文 4：完整示范——从素材到成文全流程（含批注）}

\textbf{题目}：阅读以下材料，根据要求写作。（60 分）

「有人说，这是一个『快时代』——信息秒达、交通便捷、生活节奏越来越快。但也有人说，真正的成长恰恰需要『慢下来』。」

请结合材料写一篇文章，体现你的感悟与思考。

\begin{center}
\textbf{\Large《快慢之间见真章》}
\end{center}

\textbf{【第一步：审题——圈关键词】}

快时代、慢下来、成长——这是一个思辨关系题，核心是「快与慢的辩证统一」。

\textbf{【第二步：确立中心论点】}

\boxed{\text{快是时代的节奏，慢是成长的姿态。真正的智慧，不是在快慢之间二选一，而是在快中保持慢的定力，在慢中积蓄快的能量。}}

\textbf{【第三步：安排结构 + 选择素材】}

\begin{itemize}
  \item 分论点 1（个人——因果分析）：快时代更需要慢的沉淀。\textrightarrow 素材：屠呦呦 190 次失败（慢工出细活）
  \item 分论点 2（社会——对比分析）：快不是目的，方向比速度更重要。\textrightarrow 素材：苏轼被贬（慢下来反而活得更精彩）
  \item 分论点 3（时代——升华分析）：青年当在快时代中保持慢的定力。\textrightarrow 素材：黄文秀/航天青年（落脚青年奉献）
\end{itemize}

\textbf{【第四步：写作 + 完整批注】}

\begin{center}
\textbf{完整文章}
\end{center}

「这是一个『快』的时代：信息以光速传播，飞机让世界变平，AI 在秒级之内完成人类需要数小时的工作。」（现象引入）然而，真正的成长从来不是加速度的函数——有时候，慢下来，恰恰是为了更好地快起来。」（分析过渡）\textbf{快是时代的节奏，慢是成长的姿态。在快与慢的辩证中找准自己的步调，方能在洪流中不迷失、在沉淀中真成长。}（中心论点）

\textbf{【批注：开头 150 字——现象 + 过渡 + 中心论点，出现核心关键词】}

成长需要慢的耐心。真正的突破，往往建立在对细节的极致打磨之上。\textbf{（P——分论点 1）}屠呦呦在发现青蒿素的过程中，带领团队经历了 190 次失败。她一次次翻阅古籍、调整方案，在最绝望的时候从葛洪的《肘后备急方》中获得灵感——「青蒿一握，以水二升渍，绞取汁」。这看似「慢」的反复试错，恰恰是科学最诚实的方法。\textbf{（E——叙事）}试想，如果她在第 10 次失败后就放弃了，或者为了追求速度草率宣布成果，青蒿素或许永远不会被发现。正是这种「慢」的坚持，让她在 190 次失败后迎来了那一次改变世界的成功。\textbf{（E——假设分析）}屠呦呦的故事告诉我们：有些路快不得，有些功夫省不得。\textbf{（回扣）}

\textbf{【批注：本段 230 字。P（20 字）→ E 叙事（80 字）→ E 假设分析（90 字）→ 回扣（20 字）。分析类型：假设分析——「试想，如果……」】}

快不是目的，方向比速度更重要。如果失去了方向，跑得再快也只是南辕北辙。\textbf{（P——分论点 2）}苏轼一生三度被贬，从黄州到惠州再到儋州，一次比一次偏远。在别人看来，他的仕途已经「慢」到了几乎停滞。但他没有在焦虑中追赶什么，反而在每一个「慢下来」的地方活出了最丰盛的样子：在黄州写下「大江东去」，在惠州发明荔枝酒，在儋州开办学堂。\textbf{（E——叙事）}与那些在官场中追名逐利的人不同，苏轼的可贵在于：他从不把人生的低谷当作停滞——他把每一个低谷都变成了另一种意义上的高峰。因为他心中装着一个足够大的世界，所以外在的快慢根本定义不了他。\textbf{（E——对比分析）}人生重要的不是跑得多快，而是有没有一个值得奔赴的方向。\textbf{（回扣）}

\textbf{【批注：本段 260 字。P（20 字）→ E 叙事（80 字）→ E 对比分析（120 字）→ 回扣（20 字）。分析类型：对比分析——「与……不同，……」】}

那么，身处快时代的青年，应该如何自处？\textbf{（过渡句，衔接分论点 2 和 3）}我认为，答案不是拒绝快，而是在快节奏中保持一份慢的定力——这是一种选择，更是一种担当。\textbf{（P——分论点 3）}黄文秀从北师大毕业后，没有急着挤进城市的快车道，而是选择回到广西百坭村。她用了整整一年的时间，走遍每一户人家，汽车里程正好两万五千里。\textbf{（E——叙事）}她的选择看似「慢」——放弃了城市的高薪和机遇，去做一件见效慢、甚至看不见尽头的事。但正是这种「慢」，让她真正走进了群众的心里，让百坭村在脱贫路上没有落下一户一人。\textbf{（E——因果分析）}她的事迹启示我们：\textbf{当代青年在快时代中，需要的不是焦虑地追赶，而是清醒地选择——选择那些值得慢下来深耕的领域，把个人的脚步踏在祖国最需要的地方。}这本身就是一种最深沉的爱国。\textbf{（升华分析 + 落脚青年奉献）}

\textbf{【批注：本段 280 字（重点段）。P（30 字）→ 过渡（20 字）+ E 叙事（60 字）→ E 因果分析（60 字）→ 升华分析（80 字）→ 落脚青年奉献。分析类型：因果分析 + 升华分析——从个人选择上升到家国情怀】}

快与慢从来不是对立的。快，是时代给予我们的机遇；慢，是我们回赠时代的诚意。\textbf{（总结）}作为新时代青年，我们既要有策马扬鞭、只争朝夕的干劲，也要有沉淀内心、深耕厚植的定力。在快慢之间找准自己的节奏，用慢的积累兑换快的飞跃——这，才是成长的真谛，也是我们这一代人书写青春答卷的最好方式。\textbf{（升华 + 呼吁）}

\textbf{【批注：结尾 120 字——总结快慢关系 + 上升到青年成长 + 呼吁式结尾】}

\textbf{【全文 850 字，符合 800 字要求】}

\textbf{【素材使用总结】}只用 3 个素材（屠呦呦、苏轼、黄文秀），每个都用出细节+分析，不堆砌、不空洞。

\textbf{【结构总结】}
\begin{itemize}
  \item 开头：现象引入 + 分析过渡 + 中心论点（150 字）
  \item 分论点 1（230 字）：PEE + 假设分析——屠呦呦
  \item 分论点 2（260 字）：PEE + 对比分析——苏轼
  \item 分论点 3（280 字）：PEE + 因果分析 + 升华分析——黄文秀 $\to$ 青年奉献
  \item 结尾（120 字）：总结 + 升华 + 呼吁
\end{itemize}

\textbf{【为什么这篇可以稳拿 48+？】}
\begin{enumerate}
  \item 审题精准：紧扣「快」「慢」两个关键词，不偏不倚
  \item 结构清晰：递进式（个人$\to$社会$\to$时代），逻辑链条完整
  \item 素材精当：3 个素材各从不同角度切入，不重复
  \item 论证充分：每段都有分析句（假设/对比/因果/升华），不空洞
  \item 结尾升华：自然落到青年奉献，不生硬
  \item 语言流畅：有金句、有排比、有对比，但不过度追求辞藻
\end{enumerate}
